\chapter{Mobility aids}

\section*{Definition and Overview}
Mobility aids are devices designed to assist individuals who experience difficulty walking or moving independently. These devices improve stability, reduce the weight-bearing load on lower limbs, facilitate safety, and enhance independence and participation in daily activities. Mobility aids range from simple walking sticks and crutches to walking frames (rollators) and manual or powered wheelchairs.

\section*{Epidemiology}
The use of mobility aids is highly prevalent, particularly among older adults, individuals with chronic neurological or musculoskeletal conditions, and patients recovering from surgery or trauma. In Australia, mobility impairment is one of the most common forms of disability, affecting millions of people. The demand for mobility devices is rising steadily due to an ageing population and the increasing prevalence of age-related conditions like osteoarthritis, osteoporosis, and stroke.

\section*{Aetiology and Pathophysiology}
Mobility impairment arises from a variety of physiological and anatomical deficits, which mobility aids help compensate for:

\begin{itemize}
    \item \textbf{Biomechanical instability:} Muscle weakness, joint degeneration (e.g., severe knee or hip osteoarthritis), or skeletal deformities reduce the body's ability to support itself.
    \item \textbf{Sensory and motor deficits:} Damage to the nervous system (e.g., from stroke, Parkinson's disease, or peripheral neuropathy) impairs balance, coordination, and gait.
    \item \textbf{Pain and weight-bearing limits:} Acute injuries (like fractures or sprains) or post-operative restrictions require temporary offloading of weight to allow tissue healing.
    \item \textbf{Cardiorespiratory limitation:} Severe heart or lung disease (e.g., COPD, heart failure) limits endurance, where wheeled walkers with seats provide opportunities for rest.
\end{itemize}

\section*{Clinical Features}
Individuals who require mobility aids typically present with gait and balance abnormalities:

\begin{itemize}
    \item \textbf{Gait abnormalities:} Antalgic gait (limping due to pain), broad-based gait (due to instability), or shuffling gait (characteristic of Parkinson's disease).
    \item \textbf{Balance impairment:} Frequent swaying, difficulty turning, or a history of falls or near-falls.
    \item \textbf{Reduced exercise tolerance:} Shortness of breath or severe fatigue when walking short distances.
    \item \textbf{Pain on movement:} Severe joint or muscle pain during walking or standing, which improves with rest or weight support.
\end{itemize}

\section*{Differential Diagnosis}
Before prescribing a mobility aid, the underlying cause of mobility impairment must be diagnosed and investigated:

\begin{itemize}
    \item \textbf{Reversible musculoskeletal injuries:} Soft tissue strains or fractures that require temporary rather than permanent mobility support.
    \item \textbf{Treatable neurological conditions:} Vitamin B12 deficiency, normal pressure hydrocephalus, or drug-induced parkinsonism, which can cause gait changes that may improve with treatment.
    \item \textbf{Vestibular dysfunction:} Inner ear disorders causing vertigo and imbalance, which may be treated with vestibular rehabilitation.
    \item \textbf{Severe sarcopenia or frailty:} Deconditioning that may improve with targeted strength and balance training rather than relying solely on aids.
\end{itemize}

\section*{When to Seek Medical Care}
Individuals experiencing progressive difficulty walking, frequent falls, or joint pain that limits movement should consult a GP or occupational therapist/physiotherapist for a professional assessment. Immediate medical care is needed if mobility loss is sudden or associated with trauma.

\textbf{Call 000 (or local emergency services) for:}
\begin{itemize}
    \item Sudden loss of ability to walk or move one side of the body (signs of a stroke)
    \item A fall resulting in severe pain, inability to bear weight, or a suspected limb fracture or head injury
    \item Sudden onset of severe back pain accompanied by weakness or numbness in both legs
    \item Chest pain or severe shortness of breath when trying to walk or stand
    \item Sudden loss of consciousness or severe dizziness
\end{itemize}

\section*{Diagnosis and Investigations}
A professional evaluation is necessary to select and fit the correct mobility aid:

\begin{itemize}
    \item \textbf{Physiotherapist assessment:} Detailed gait analysis, balance testing, and upper and lower limb strength assessment.
    \item \textbf{Occupational therapy assessment:} Evaluating the patient's home environment, activities of daily living (ADLs), and cognitive ability to use the device safely.
    \item \textbf{Radiological imaging:} X-rays, CT, or MRI scans of the spine, hips, or knees to diagnose underlying joint or bone pathology.
    \item \textbf{Cardiovascular and neurological workup:} Electrocardiogram (ECG), echocardiogram, or brain imaging if transient ischaemic attacks (TIAs) or cardiac syncope is suspected.
\end{itemize}

\section*{Management}
The management of mobility issues involves matching the patient with the correct device, training, and ongoing rehabilitation:

\begin{itemize}
    \item \textbf{Device selection:} Choosing the most appropriate aid, such as single-point canes for mild balance issues, quad canes for extra stability, underarm or forearm crutches for temporary offloading, rollators (walkers with wheels, brakes, and seats) for endurance, or wheelchairs for severe mobility loss.
    \item \textbf{Fitting and adjustment:} Ensuring the device is at the correct height (typically matching the crease of the wrist when standing straight) to prevent musculoskeletal strain.
    \item \textbf{Gait training:} Instruction by a physiotherapist on correct usage patterns, such as placing the cane on the side opposite the weak leg.
    \item \textbf{Environmental modifications:} Removing trip hazards (like loose rugs), installing grab rails, and widening doorways to facilitate safe device use at home.
    \item \textbf{Physical rehabilitation:} Ongoing strength and balance exercises (such as the Otago Exercise Programme) to improve natural mobility and reduce device dependence where possible.
\end{itemize}

\section*{Complications}
\begin{itemize}
    \item Fall injuries due to incorrect use of the device, poor maintenance (e.g., worn rubber tips), or trip hazards
    \item Postural and musculoskeletal pain (e.g., shoulder, wrist, or back strain) from an poorly fitted or inappropriately sized aid
    \item Nerve compression (e.g., crutch palsy affecting the radial nerve) from incorrect crutch usage
    \item Muscle deconditioning and joint stiffness due to over-reliance on a mobility aid when active exercise is possible
    \item Pressure injuries or skin breakdown from prolonged wheelchair use
    \item Social isolation or psychological distress related to stigma or frustration with mobility device dependence
\end{itemize}

\section*{Prognosis and Outcomes}
The use of a well-fitted mobility aid significantly improves the user's safety, confidence, and quality of life. For patients recovering from acute injuries or surgery, mobility aids are temporary tools that facilitate safe mobilisation during the healing process. For those with chronic, progressive conditions, these devices are essential lifelines that prevent falls, prolong independent living, and allow continued participation in social and community activities.

\section*{Prevention and Self-Care}
\begin{itemize}
    \item Inspect your mobility aid weekly for loose parts, cracked frames, or worn rubber tips, and replace them promptly
    \item Keep floors clear of clutter, electrical cords, and loose rugs, and ensure all walking areas are well-lit
    \item Wear supportive, flat, non-slip footwear when using your mobility device; avoid walking in socks or loose slippers
    \item Learn how to navigate stairs and uneven terrain safely with your aid under the guidance of a physiotherapist
    \item Store your mobility aid within easy reach of your bed or chair to prevent falls when getting up
\end{itemize}

\section*{Sources and Further Reading}
The following organisations provide reliable information on mobility aids:

\begin{itemize}
    \item Independent Living Assessment (Australia). \textit{Guides on choosing assistive technology and mobility devices.} https://ilaportal.org.au/
    \item Healthdirect Australia. \textit{Information on mobility aids, selection, and safety tips.} https://www.healthdirect.gov.au/
    \item Mayo Clinic. \textit{Patient guide to choosing and using canes and walkers.} https://www.mayoclinic.org/
    \item NHS. \textit{Guide to walking aids, wheelchairs, and getting a professional assessment.} https://www.nhs.uk/
    \item Australian Physiotherapy Association. \textit{Information on gait training, falls prevention, and rehabilitation.} https://choose.physio/
\end{itemize}
